\chapter{Firebase Database Design and Management}

\section{Introduction}

This chapter describes the design, structure, and management of the database used in the ENDOTECH system. The system utilizes Firebase Cloud Firestore as the primary database solution, enabling real-time data synchronization and scalable cloud storage.

The database plays a critical role in ensuring data consistency, availability, and efficient processing of inventory and sales operations.

\section{Firebase Overview and Selection Rationale}

Firebase is a Backend-as-a-Service (BaaS) platform developed by Google. It provides a set of cloud-based tools for application development, including authentication, hosting, and real-time databases.

Cloud Firestore was selected as the database for the ENDOTECH system due to the following advantages:

\begin{itemize}
    \item Real-time data synchronization across multiple devices;
    \item High scalability without manual server management;
    \item Flexible NoSQL document-based data model;
    \item Built-in authentication and security rules;
    \item Offline support for mobile applications.
\end{itemize}

These features make Firebase particularly suitable for cross-platform applications developed with Flutter.

\section{Database Architecture}

The database follows a NoSQL document-oriented structure. Data is organized into collections and documents rather than tables and rows.

Each collection contains multiple documents, and each document consists of key-value pairs.

The main collections in the ENDOTECH system include:

\begin{itemize}
    \item \texttt{users}
    \item \texttt{inventory}
    \item \texttt{sales}
    \item \texttt{categories}
    \item \texttt{suppliers}
\end{itemize}

This structure allows flexible and scalable data storage.

\section{Data Model and Structure}

\subsection{Users Collection}

The \texttt{users} collection stores information about system users.

\begin{lstlisting}[language=JSON, caption=User Document Structure]
{
  "userId": "string",
  "email": "string",
  "displayName": "string",
  "role": "admin | manager | staff",
  "createdAt": "timestamp",
  "lastLogin": "timestamp"
}
\end{lstlisting}

\subsection{Inventory Collection}

The \texttt{inventory} collection contains medical equipment data.

\begin{lstlisting}[language=JSON, caption=Inventory Document Structure]
{
  "itemId": "string",
  "name": "string",
  "description": "string",
  "categoryId": "reference",
  "supplierId": "reference",
  "quantity": "number",
  "price": "number",
  "minStockLevel": "number",
  "createdAt": "timestamp",
  "updatedAt": "timestamp"
}
\end{lstlisting}

\subsection{Sales Collection}

The \texttt{sales} collection records transactions.

\begin{lstlisting}[language=JSON, caption=Sales Document Structure]
{
  "saleId": "string",
  "userId": "reference",
  "items": [
    {
      "itemId": "reference",
      "quantity": "number",
      "price": "number"
    }
  ],
  "totalAmount": "number",
  "timestamp": "timestamp"
}
\end{lstlisting}

\section{Data Relationships and Logical Model}

Although Firestore is a NoSQL database, logical relationships are maintained through references.

\begin{itemize}
    \item Each sale is linked to a user;
    \item Sales contain references to inventory items;
    \item Inventory items are associated with categories and suppliers;
    \item Users are assigned roles for access control.
\end{itemize}

This approach ensures flexibility while maintaining logical consistency.

\section{Database Operations}

\subsection{Reading Data}

\begin{lstlisting}[language=Dart]
FirebaseFirestore.instance
  .collection('inventory')
  .orderBy('name')
  .snapshots();
\end{lstlisting}

\subsection{Writing Data}

\begin{lstlisting}[language=Dart]
FirebaseFirestore.instance.collection('sales').add({
  'userId': FirebaseAuth.instance.currentUser!.uid,
  'totalAmount': total,
  'timestamp': FieldValue.serverTimestamp(),
});
\end{lstlisting}

\subsection{Real-Time Updates}

Firestore provides real-time listeners:

\begin{lstlisting}[language=Dart]
FirebaseFirestore.instance
  .collection('inventory')
  .snapshots();
\end{lstlisting}

\section{Security and Access Control}

Security rules ensure that only authorized users can access data.

\begin{lstlisting}[language=JavaScript]
rules_version = '2';
service cloud.firestore {
  match /databases/{database}/documents {
    match /inventory/{itemId} {
      allow read, write: if request.auth != null;
    }
  }
}
\end{lstlisting}

These rules restrict access to authenticated users.

\section{Performance Optimization}

To improve performance, the following strategies are used:

\begin{itemize}
    \item Indexing frequently queried fields;
    \item Minimizing data transfer;
    \item Using real-time listeners efficiently;
    \item Implementing local caching.
\end{itemize}

\section{Backup and Data Management}

Firebase provides tools for data backup and export. Regular backups can be performed using Firebase CLI and Google Cloud services.

This ensures data safety and recovery in case of failures.

\section{Conclusion}

The database design of the ENDOTECH system provides a scalable, flexible, and efficient solution for managing inventory and sales data. The use of Firebase Cloud Firestore enables real-time synchronization, secure data access, and seamless integration with the Flutter application.